\subsection{ASR Experimental Methodology}
\label{sec:asr_methodology}

\paragraph{Hardware and software.}
Experiments were conducted on an NVIDIA H200 MIG
\texttt{1g.18gb} GPU partition with 16 streaming multiprocessors
and compute capability 9.0 (\texttt{sm\_90}). The final matched
end-to-end comparisons used PyTorch 2.10.0 (CUDA 12.8 build) with CUDA 13.2.
The custom CUDA library was compiled for \texttt{sm\_90}.

\paragraph{Kernel-level evaluation.}
We evaluated GEMM tile sizes $T\in\{8,16,32\}$ using the
supplied Task 3 benchmark. Each configuration was checked
against the reference implementation. We compared latency
and achieved floating-point throughput against the supplied
naive GEMM and cuBLAS baselines across encoder, decoder,
and language-model-head matrix shapes.

\paragraph{End-to-end evaluation.}
The ASR benchmark processed the same 3.50-second audio
sample, represented by mel features of shape
$(1,128,3000)$, with a 59-token prompt and 13 generated
tokens. Each timed configuration used one warm-up run
followed by three measured runs. We report mean
transcription latency, peak GPU memory, and the
supplied transcription correctness check.

Four configurations were compared: (i) PyTorch-only
operations; (ii) all supported custom CUDA kernels;
(iii) custom CUDA except GEMM, which used PyTorch;
and (iv) custom CUDA except GEMM and attention,
which both used PyTorch. All matched end-to-end
measurements used the same GPU and software environment.

\paragraph{Profiling and interpretation.}
Operation-level profiling was performed separately for
the fully custom and fastest hybrid configurations.
We report operation call counts and cumulative measured
GPU time. Profiling results are used to identify
bottlenecks, while controlled kernel substitutions
provide stronger evidence about the effect of
implementation choices on end-to-end latency.

The correctness check covers one supplied transcription
sample and is not a corpus-level word error rate.
Similarly, the reported latency results characterise
the tested hardware, model and workload rather than
all possible ASR inference settings.
